% Generated draft: requires mathematical and source review.
\documentclass[11pt]{article}
\usepackage[margin=25mm]{geometry}
\usepackage{fontspec}
\setmainfont{Latin Modern Roman}
\usepackage{amsmath,amssymb,mathtools,unicode-math}
\setmathfont{Latin Modern Math}
\usepackage{enumitem,needspace,xurl,hyperref}
\setlength{\parindent}{0pt}
\setlength{\parskip}{5pt}
\setlength{\emergencystretch}{3em}
\newcommand{\N}{\mathbb{N}}
\newcommand{\Z}{\mathbb{Z}}
\newcommand{\Q}{\mathbb{Q}}
\newcommand{\R}{\mathbb{R}}
\newcommand{\C}{\mathbb{C}}
\newcommand{\F}{\mathbb{F}}
\newcommand{\A}{\mathbb{A}}
\newcommand{\PP}{\mathbb{P}}
\newcommand{\E}{\mathbb{E}}
\newcommand{\eps}{\varepsilon}
\newcommand{\dd}{\,\mathrm{d}}
\begin{document}
\section*{UnsolvedMath ID: 30004444}
\textbf{Source problem code:} catalogue entry
\subsection{Definitions and mathematical statement}
Work in \(\mathsf{ZFC}\). A cardinal \(\kappa\) is singular if \(\operatorname{cf}(\kappa)<\kappa\), where \(\operatorname{cf}(\alpha)\) is the least order type of a cofinal subset of the ordinal \(\alpha\). It is a strong-limit cardinal if \(2^\lambda<\kappa\) for every cardinal \(\lambda<\kappa\). Write \(\kappa^+\) for the successor cardinal of \(\kappa\), and \(2^\kappa\) for the cardinality of its power set.

The following conventional definitions make the displayed questions explicit; the supplied record does not establish that these are precisely the conventions of the original report. For singular \(\kappa\), put \(\delta=\operatorname{cf}(\kappa)\). Its pseudopower \(\operatorname{pp}(\kappa)\) is conventionally defined by
\[
\operatorname{pp}(\kappa)=\sup\left\{\operatorname{tcf}\left(\prod_{i<\delta}\lambda_i/J\right):\begin{array}{l}
\lambda_i<\kappa\text{ are infinite regular cardinals},\\
J\text{ is a proper ideal on }\delta\text{ containing all bounded subsets},\\
\{i<\delta:\lambda_i\leq\mu\}\in J\text{ for every cardinal }\mu<\kappa,\\
\operatorname{tcf}\left(\prod_{i<\delta}\lambda_i/J\right)\text{ exists}
\end{array}\right\}.
\]
Here the product consists of functions \(f\) with \(f(i)<\lambda_i\), and
\[
f<_J g\quad\Longleftrightarrow\quad\{i<\delta:f(i)\geq g(i)\}\in J.
\]
Its true cofinality \(\operatorname{tcf}\) is the regular cardinal length of a strictly \(<_J\)-increasing cofinal sequence, when such a sequence exists.

Under the conventional interpretation of weak square at \(\kappa\) as \(\square_\kappa^*\), this principle asserts the existence of a sequence
\[
\langle\mathcal C_\alpha:\alpha<\kappa^+,\ \alpha\text{ a nonzero limit ordinal}\rangle
\]
such that each \(\mathcal C_\alpha\) is a nonempty family of at most \(\kappa\) clubs in \(\alpha\), every member has order type at most \(\kappa\), and whenever \(C\in\mathcal C_\alpha\) and \(\beta\) is a limit point of \(C\), one has \(C\cap\beta\in\mathcal C_\beta\). A club is a closed unbounded subset in the ordinal order topology. The intended weak-square convention requires confirmation from the report.

For a regular uncountable cardinal \(\theta\), a \(\theta\)-tree is a partially ordered set whose predecessors below each node are well ordered, whose height is \(\theta\), and whose levels all have cardinality less than \(\theta\). The tree property \(\operatorname{TP}(\theta)\) asserts that every such tree has a branch meeting cofinally many levels.

\begin{enumerate}
\item The displayed question attributed to Shelah by the supplied assessment asks whether
\[
\forall\kappa\;\bigl(\kappa\text{ is a singular cardinal}\ \Longrightarrow\ \operatorname{cf}(\operatorname{pp}(\kappa))>\kappa\bigr).
\]
\item The supplied statement also asks whether either of the following configurations can occur:
\[
\exists\kappa\;\bigl(\kappa\text{ is a singular strong-limit cardinal}\ \land\ \neg\square_\kappa^*\ \land\ 2^\kappa>\kappa^+\bigr),
\]
or
\[
\exists\kappa\;\bigl(\kappa\text{ is a singular strong-limit cardinal}\ \land\ \operatorname{TP}(\kappa^+)\ \land\ 2^\kappa>\kappa^+\bigr).
\]
These are presented as separate existence or consistency questions; the text does not require the same \(\kappa\) to witness both configurations. No large-cardinal assumptions or intended consistency-strength target are supplied.
\end{enumerate}

There is a material qualification: although the supplied statement explicitly lists both configurations, the accompanying assessment says that the extraction truncates immediately after introducing a separate Woodin question. Thus the displayed formulas can be retained as supplied questions, but the completeness of the intended problem and the attribution of the later questions remain unresolved. The supplied dated assessment describes the first question as answered negatively by extender-based forcing and retains an overall open label because the extraction is incomplete. These are supplied curation claims, not independently verified mathematical conclusions.
\subsection{Short English statement}
The record asks whether the pseudopower of every singular cardinal has cofinality larger than that cardinal, and whether a singular strong-limit cardinal can violate the singular cardinal hypothesis while weak square fails or the tree property holds at its successor. The supplied assessment reports a negative answer to the first question but also reports an incomplete extraction involving a separate Woodin question. The remaining scope, conventions, and attribution require review.
\subsection{Sources}
\begin{itemize}
\item Local registry: \texttt{data/unsolvedmath/problems.json}, with \texttt{catalog/problems.json} as the additional-record registry; supplied record for permanent UnsolvedMath ID \(30004444\), source code \texttt{catalogue entry}, title \emph{Singular Cardinal Arithmetic Without Weak Square}. Registry contents have not been independently inspected for this draft.
\item Supplied source citation: \emph{Set Theory (2021)}, Oberwolfach Reports, pp.~797--855, DOI \texttt{10.4171/owr/2020/14}. \url{https://doi.org/10.4171/owr/2020/14}. The background also labels this source \emph{Oberwolfach Report 14/2020}; the record gives proposed year 2020.
\item Supplied literature reference: Gitik, \emph{Extender-based forcings} (2020). \url{https://doi.org/10.1142/S0219061320500130}.
\item Supplied literature reference: \emph{Oberwolfach Report 14/2020}. \url{https://doi.org/10.4171/owr/2020/14}.
\item Literature checked at: \texttt{2026-08-22}. This exact date records supplied dated curation, not independent verification. The supplied assessment reports that extender-based forcing answers the displayed Shelah question negatively, but says the extraction is incomplete after introducing a separate Woodin question and therefore assigns the full record an open status. The background explicitly calls for independent verification of time-sensitive status and cited claims.
\item Attribution is limited to the assessment's association of the first question with Shelah and its mention of a separate Woodin question. The structured proposer field is empty. The supplied \texttt{corrected\_verified} statement label is qualified as repairs only to delimiters, notation, grammar, and typography, preserving mathematical content and scope.
\end{itemize}
\end{document}
